\section{Preprocessing, evaluation, and simple baselines}
\label{sec:preprocessing-evaluation-simple-baselines}

\subsection{Purpose of this stage}

This section starts the modelling workflow after the raw-data audit, deterministic correction, train/test split, and training-set exploratory analysis. The purpose is not yet to fit flexible learned models. Instead, this stage establishes the evaluation language and simple reference baselines that all later models must improve upon.

The held-out test set remains unused. All results in this section are computed only from the training set using stratified cross-validation. This preserves the final test set for the final evaluation stage.

The section has three goals. First, it defines the binary classification evaluation framework used throughout the rest of the project. Second, it introduces reusable preprocessing infrastructure for later models. Third, it evaluates simple baselines that do not learn a flexible relationship from the full feature matrix.

The simple baselines considered here are:
\begin{enumerate}
    \item a majority-class baseline;
    \item a prior-probability baseline;
    \item a stratified random baseline;
    \item a uniform random baseline;
    \item an exploratory-analysis-inspired rule baseline.
\end{enumerate}

Learned models such as least-squares classification, logistic regression, k-nearest neighbours, Naive Bayes, decision trees, support vector machines, and neural networks are introduced in later sections.

\subsection{Training-only evaluation discipline}

Let the training set be
\[
\mathcal{D}_{train}
=
\{(x_i,y_i)\}_{i=1}^{n},
\]
where \(x_i\) denotes the feature vector for customer \(i\), and \(y_i \in \{0,1\}\) denotes the binary churn label. In this project, the positive class is churn:
\[
y_i = 1 \quad \Longleftrightarrow \quad \text{customer } i \text{ churned}.
\]

All model-development decisions must be made without using the held-out test set. This includes preprocessing design, feature engineering, feature selection, model-family choice, hyperparameter tuning, threshold tuning, calibration decisions, and model comparison.

For this stage, performance is estimated using stratified 5-fold cross-validation inside the training set. Stratification preserves approximately the same churn rate in each fold. This is important because churn is the minority class. In the training set, the class distribution is shown in Table~\ref{tab:stage04-target-distribution}.

\begin{table}[H]
\centering
\caption{Training-set target distribution.}
\label{tab:stage04-target-distribution}
\begin{tabular}{lrr}
\toprule
Class & Count & Percentage \\
\midrule
No churn (\(Y=0\)) & 4139 & 73.46\% \\
Churn (\(Y=1\)) & 1495 & 26.54\% \\
\bottomrule
\end{tabular}
\end{table}

Cross-validation produces out-of-fold predictions. This means that the prediction for each observation is generated by a model that was not fitted on that same observation. This gives a more honest development-stage estimate than evaluating a model on the same data used to fit it.

\subsection{Positive and negative classes}

Many binary classification metrics are defined relative to the positive class. The positive class is not necessarily the most common class. It is the event of interest.

A medical screening example makes the terminology intuitive. Suppose a test predicts whether a person has cancer. Then:
\begin{itemize}
    \item positive class: the person has cancer;
    \item negative class: the person does not have cancer;
    \item predicted positive: the test says cancer;
    \item predicted negative: the test says no cancer.
\end{itemize}

For the churn problem, the same logic becomes:
\begin{itemize}
    \item positive class: the customer churns;
    \item negative class: the customer does not churn;
    \item predicted positive: the model flags the customer as likely to churn;
    \item predicted negative: the model does not flag the customer as likely to churn.
\end{itemize}

This convention matters because recall, precision, \(F_1\), ROC-AUC, and PR-AUC all depend on which class is treated as positive.

\subsection{Confusion matrix}

The confusion matrix counts the four possible combinations of actual class and predicted class. The positive-first layout used here is shown in Table~\ref{tab:stage04-generic-confusion-matrix}. Rows represent the actual class and columns represent the predicted class.

\begin{table}[H]
\centering
\caption{Positive-first binary confusion matrix layout.}
\label{tab:stage04-generic-confusion-matrix}
\begin{tabular}{llcc}
\toprule
 & & \multicolumn{2}{c}{Predicted} \\
 & & Positive & Negative \\
\midrule
Actual & Positive & TP & FN \\
Actual & Negative & FP & TN \\
\bottomrule
\end{tabular}
\end{table}

The entries are:
\[
TP = \#\{\hat{y}=1, y=1\},
\qquad
FN = \#\{\hat{y}=0, y=1\},
\]
\[
FP = \#\{\hat{y}=1, y=0\},
\qquad
TN = \#\{\hat{y}=0, y=0\}.
\]

This order is chosen to match the lecture-style visual explanation: actual positives are shown first, and within that row the model can either correctly predict positive or incorrectly predict negative. Some software libraries store confusion matrices in a negative-first numeric order. In the code, the counts are computed with explicit labels and then renamed so that the report can present the matrix in the positive-first order.

In the medical screening example:
\begin{itemize}
    \item true positive: the test says cancer, and the person really has cancer;
    \item false negative: the test says no cancer, but the person really has cancer;
    \item false positive: the test says cancer, but the person does not have cancer;
    \item true negative: the test says no cancer, and the person does not have cancer.
\end{itemize}

In the churn problem:
\begin{itemize}
    \item true positive: the model predicts churn, and the customer churns;
    \item false negative: the model predicts no churn, but the customer churns;
    \item false positive: the model predicts churn, but the customer does not churn;
    \item true negative: the model predicts no churn, and the customer does not churn.
\end{itemize}

False positives and false negatives have different practical meanings. A false positive may lead to unnecessary retention effort for a customer who would have stayed anyway. A false negative may mean missing a customer who actually leaves.

\subsection{Accuracy and the class-imbalance problem}

Accuracy is the fraction of all predictions that are correct:
\[
\text{Accuracy}
=
\frac{TP + TN}{TP + FN + FP + TN}.
\]

Accuracy is intuitive, but it can be misleading when one class is much more common than the other. In this dataset, about \(73.46\%\) of training customers do not churn. A classifier that always predicts no churn can therefore achieve \(73.46\%\) accuracy while detecting zero churners.

This is the central reason why the project does not evaluate churn models using accuracy alone. The evaluation must also measure how well the model detects the positive class and how many false churn alerts it creates.

\subsection{Recall, specificity, precision, and \(F_1\)}

Recall, also called true positive rate, measures the fraction of actual positives that are detected:
\[
\text{Recall}
=
\text{TPR}
=
\frac{TP}{TP+FN}.
\]
In this project, recall answers: among all customers who churned, how many did the model flag?

Specificity, also called true negative rate, measures the fraction of actual negatives correctly rejected:
\[
\text{Specificity}
=
\text{TNR}
=
\frac{TN}{TN+FP}.
\]
In this project, specificity answers: among all customers who did not churn, how many did the model correctly leave unflagged?

Precision measures the reliability of positive predictions:
\[
\text{Precision}
=
\frac{TP}{TP+FP}.
\]
In this project, precision answers: among all customers flagged as churners, how many actually churned?

The \(F_1\)-score combines precision and recall:
\[
F_1
=
2\cdot
\frac{\text{Precision}\cdot\text{Recall}}
{\text{Precision}+\text{Recall}}.
\]
It is useful as a compact summary when both precision and recall matter. However, it hides the separate business meanings of false positives and false negatives. For this reason, \(F_1\) is reported together with the underlying precision and recall values.

\subsection{Balanced accuracy}

Balanced accuracy averages recall and specificity:
\[
\text{Balanced Accuracy}
=
\frac{1}{2}
\left(
\text{Recall}+\text{Specificity}
\right).
\]

Balanced accuracy is useful under class imbalance because it gives equal importance to the positive and negative classes. A classifier that always predicts one class receives balanced accuracy equal to \(0.5\) in a binary problem, because it perfectly classifies one class and completely fails on the other.

This makes balanced accuracy useful for comparing simple baselines. A majority-class classifier may have high ordinary accuracy, but its balanced accuracy reveals whether it is actually useful for both classes.

\subsection{ROC-AUC and PR-AUC}

Many classifiers produce a score or probability rather than only a hard class label. A threshold turns this score into a class prediction:
\[
\hat{y}_{\tau}
=
\mathbb{1}
\{
s(x)\geq \tau
\},
\]
where \(s(x)\) is a model score or predicted probability and \(\tau\) is the decision threshold.

Changing the threshold changes the tradeoff between false positives and false negatives. ROC and precision-recall curves summarize this threshold-dependent behaviour.

The ROC curve plots:
\[
\text{FPR}
=
\frac{FP}{FP+TN}
\quad
\text{against}
\quad
\text{TPR}
=
\frac{TP}{TP+FN}.
\]
ROC-AUC summarizes how well the model ranks positive cases above negative cases across thresholds.

The precision-recall curve plots:
\[
\text{Precision}
=
\frac{TP}{TP+FP}
\quad
\text{against}
\quad
\text{Recall}
=
\frac{TP}{TP+FN}.
\]
PR-AUC summarizes how precise the positive predictions remain as the model attempts to recover more positives.

ROC-AUC and PR-AUC often move in the same direction when one model is clearly better at ranking positive cases. They can differ under class imbalance. A model can have a reasonable ROC-AUC while still having low precision if many of its predicted positives are false positives. Since churn is the minority class and the positive class is the main event of interest, PR-AUC is an important complement to ROC-AUC.

\subsection{Thresholds and calibration}

A probabilistic classifier estimates:
\[
\hat{p}(Y=1\mid X=x).
\]
The default hard classification rule is often:
\[
\hat{y}
=
\mathbb{1}
\{
\hat{p}(Y=1\mid X=x)\geq 0.5
\}.
\]

The threshold \(0.5\) is not automatically optimal. In churn prediction, the preferred threshold depends on the relative cost of offering retention actions to customers who would not churn versus missing customers who will churn.

Calibration is a separate question. A model is calibrated if predicted probabilities behave like empirical probabilities:
\[
P(Y=1 \mid \hat{p}=q) \approx q.
\]
For example, among customers assigned a churn probability near \(0.30\), about \(30\%\) should actually churn. Calibration will become more important when learned probability models are introduced.

\subsection{Class imbalance strategies}

The current section evaluates baselines under the natural training distribution. Resampling strategies such as random oversampling, random undersampling, and SMOTE are not applied yet because they are training interventions for learned models rather than baseline definitions.

The validation and test distributions should remain realistic. If resampling is used later, it must be applied only inside the training part of each cross-validation split. The validation fold must remain untouched. The incorrect workflow is:
\[
\text{resample full training set}
\quad \rightarrow \quad
\text{cross-validate on resampled data}.
\]
This can leak information or artificially improve validation performance because the validation fold no longer represents an untouched natural sample.

The correct workflow is:
\[
\text{split fold}
\quad \rightarrow \quad
\text{fit preprocessing on training fold}
\quad \rightarrow \quad
\text{resample training fold only}
\quad \rightarrow \quad
\text{fit model}
\quad \rightarrow \quad
\text{evaluate on untouched validation fold}.
\]
In Python, this usually requires an imbalanced-learn pipeline rather than an ordinary scikit-learn pipeline, because the sampler changes both \(X\) and \(y\). Resampling, class weighting, and threshold tuning are considered later after learned models have been introduced.

\subsection{Reusable preprocessing infrastructure}

Most simple baselines in this section do not require feature preprocessing. Dummy classifiers ignore the feature matrix, and the rule-based classifier works directly with a small number of original categorical variables.

Nevertheless, reusable preprocessing objects are introduced here because later model sections depend on them. Two broad preprocessing variants are prepared:
\begin{itemize}
    \item a scaled preprocessor for scale-sensitive models such as logistic regression, k-nearest neighbours, support vector machines, and neural networks;
    \item an unscaled preprocessor for tree-based models, where numeric scaling is usually unnecessary.
\end{itemize}

Categorical variables are one-hot encoded in standard modelling pipelines. Structural categories such as \texttt{No internet service} and \texttt{No phone service} are treated as valid categories, not as missing values.

\subsection{Simple baseline definitions}

The majority-class baseline always predicts the most frequent class observed in the training fold:
\[
\hat{y}
=
\arg\max_{c\in\{0,1\}}
\sum_{i=1}^{n}
\mathbb{1}\{y_i=c\}.
\]
In this dataset, this means always predicting no churn.

The prior-probability baseline uses the empirical class probabilities:
\[
\hat{p}(Y=c)
=
\frac{1}{n}
\sum_{i=1}^{n}
\mathbb{1}\{y_i=c\}.
\]
Its hard class prediction is still the majority class, but its predicted probabilities correspond to the class prior.

The stratified random baseline randomly samples labels according to the training-fold class distribution:
\[
\hat{Y}
\sim
\operatorname{Categorical}
\left(
\hat{p}(Y=0),
\hat{p}(Y=1)
\right).
\]

The uniform random baseline samples labels uniformly:
\[
\hat{Y}
\sim
\operatorname{Categorical}(0.5,0.5).
\]

The exploratory-analysis-inspired rule baseline assigns one risk point for each high-risk condition:
\[
s(x)
=
\sum_{j=1}^{m}
\mathbb{1}
\{
\text{condition }j\text{ is true}
\}.
\]
The predicted class is:
\[
\hat{y}
=
\mathbb{1}
\{
s(x)\geq t
\}.
\]
The rule used here sets \(t=2\), with the following high-risk conditions:
\begin{itemize}
    \item month-to-month contract;
    \item electronic check payment;
    \item fiber optic internet service;
    \item no online security;
    \item no tech support.
\end{itemize}
This rule is intentionally simple. It is not a final model. Its purpose is to test whether the strongest training-set EDA patterns already contain predictive signal.

\subsection{Simple baseline results}

Table~\ref{tab:stage04-simple-baseline-metrics} reports the cross-validated metrics for the simple baselines. Table~\ref{tab:stage04-simple-baseline-confusion} reports the corresponding confusion-matrix counts in the positive-first order \(TP,FN,FP,TN\).

\begin{table}[H]
\centering
\caption{Cross-validated simple-baseline metrics on the training set.}
\label{tab:stage04-simple-baseline-metrics}
\resizebox{\textwidth}{!}{
\begin{tabular}{lrrrrrrrrrr}
\toprule
Model & Acc. & Bal. acc. & Prec. & Recall & Spec. & \(F_1\) & ROC-AUC & PR-AUC & Pred. pos. & Obs. pos. \\
\midrule
EDA-inspired rule & 0.6076 & 0.7032 & 0.3956 & 0.9070 & 0.4994 & 0.5509 & 0.8109 & 0.5435 & 0.6084 & 0.2654 \\
Dummy: stratified random & 0.6140 & 0.5065 & 0.2748 & 0.2776 & 0.7354 & 0.2762 & 0.5065 & 0.2680 & 0.2680 & 0.2654 \\
Dummy: uniform random & 0.5051 & 0.5057 & 0.2698 & 0.5070 & 0.5045 & 0.3522 & 0.5000 & 0.2654 & 0.4986 & 0.2654 \\
Dummy: majority class & 0.7346 & 0.5000 & 0.0000 & 0.0000 & 1.0000 & 0.0000 & 0.5000 & 0.2654 & 0.0000 & 0.2654 \\
Dummy: prior probability & 0.7346 & 0.5000 & 0.0000 & 0.0000 & 1.0000 & 0.0000 & 0.4999 & 0.2653 & 0.0000 & 0.2654 \\
\bottomrule
\end{tabular}
}
\end{table}

\begin{table}[H]
\centering
\caption{Cross-validated confusion-matrix counts for the simple baselines, shown in positive-first order.}
\label{tab:stage04-simple-baseline-confusion}
\begin{tabular}{lrrrr}
\toprule
Model & TP & FN & FP & TN \\
\midrule
EDA-inspired rule & 1356 & 139 & 2072 & 2067 \\
Dummy: stratified random & 415 & 1080 & 1095 & 3044 \\
Dummy: uniform random & 758 & 737 & 2051 & 2088 \\
Dummy: majority class & 0 & 1495 & 0 & 4139 \\
Dummy: prior probability & 0 & 1495 & 0 & 4139 \\
\bottomrule
\end{tabular}
\end{table}

\subsection{Interpretation}

The majority-class and prior-probability baselines have the highest ordinary accuracy among the simple baselines, \(0.7346\). However, this accuracy is misleading. Both baselines predict no churn for every customer. Therefore, they classify all 4139 non-churners correctly but miss all 1495 churners. Their recall is exactly \(0.0000\), and their balanced accuracy is only \(0.5000\).

The stratified random baseline predicts churn at approximately the observed training-set churn rate. Its predicted positive rate is \(0.2680\), close to the observed positive rate \(0.2654\). Its balanced accuracy is \(0.5065\), which is only slightly above random-level performance. This is expected because the predictions are random and not informed by the features.

The uniform random baseline predicts churn for about half of the customers. This gives recall \(0.5070\), higher than the majority-class baseline, but this is achieved by producing many false positives. The model incorrectly flags 2051 non-churners as churners. Its ROC-AUC is \(0.5000\), as expected for a random classifier.

The EDA-inspired rule is the only simple baseline that uses feature information. It detects 1356 of the 1495 churners, giving recall \(0.9070\). This confirms that the strongest EDA patterns contain substantial predictive signal. However, it also incorrectly flags 2072 non-churners, so its specificity is only \(0.4994\) and its precision is \(0.3956\).

This makes the EDA-inspired rule a high-recall, low-specificity baseline. It is useful because it shows that simple churn-risk conditions already recover most churners. It is not sufficient as a final model because it predicts churn for about \(60.84\%\) of customers, much higher than the observed churn rate of \(26.54\%\). A practical learned model should aim to keep much of this recall while improving precision and specificity.

\subsection{Implications for later models}

The simple baselines establish the minimum standard for later learned models. A useful learned model should:
\begin{itemize}
    \item clearly outperform dummy baselines;
    \item avoid relying only on ordinary accuracy;
    \item improve balanced accuracy relative to random baselines;
    \item detect churners with non-trivial recall;
    \item improve precision and specificity relative to the broad EDA rule;
    \item provide useful ranking performance through ROC-AUC and PR-AUC;
    \item eventually support threshold tuning and calibration analysis.
\end{itemize}

The next modelling stage begins the learned-model sequence with linear classification, least-squares classification, logistic regression, and regularized logistic regression.
